% correlatividades historia
\thispagestyle{empty}
\PageDecor{\CorrPageNo}
\SectionTag{CORRELATIVIDADES}
\DocHeading{caca 24}

\begingroup
\small
\setlength{\tabcolsep}{4.5pt}
\renewcommand{\arraystretch}{1.34}
\setlength{\LTleft}{0pt}
\setlength{\LTright}{0pt}
\setlength{\LTpre}{0.35em}
\setlength{\LTpost}{0.75em}

\newcommand{\CorrTableStart}{
\begin{longtable}{|>{\RaggedRight\arraybackslash}p{0.30\textwidth}|>{\RaggedRight\arraybackslash}p{0.64\textwidth}|}
\hline
\textbf{Unidad Curricular} & \textbf{Correlativa con} \\
\hline
\endfirsthead
\hline
\textbf{Unidad Curricular} & \textbf{Correlativa con} \\
\hline
\endhead
}

\newcommand{\CorrTableEnd}{
\end{longtable}
}

\newcommand{\CorrRow}[2]{
\rule{0pt}{2.55ex}#1 & #2 \\
\hline
}

\DocQuestion{Tabla 1. Correlatividades de 2\textsuperscript{o} año}
\CorrTableStart
\CorrRow{Filosofía}{Pedagogía (R).}
\CorrRow{Psicología Educacional}{Pedagogía (R).}
\CorrRow{Práctica Docente II}{Práctica Docente I (A); Pedagogía (R); Didáctica General (A).}
\CorrRow{Sujetos de la Educación Secundaria}{Didáctica General (R); Pedagogía (R).}
\CorrRow{Procesos sociales, políticos, económicos y culturales del Feudalismo y la Modernidad}{Procesos sociales, políticos, económicos y culturales de la Antigüedad (A); Historia de las Ideas I (A); Procesos sociales, políticos, económicos y culturales de los pueblos originarios de América (R); Problemática del conocimiento histórico (R).}
\CorrRow{Procesos sociales, políticos, económicos y culturales Americanos I}{Procesos sociales, políticos, económicos y culturales de los pueblos originarios de América (R).}
\CorrRow{Historia de las ideas II}{Historia de las Ideas I (A); Problemática del conocimiento histórico (R).}
\CorrRow{Economía Política}{Procesos sociales, políticos, económicos y culturales de la Antigüedad (R); Historia de las Ideas I (R).}
\CorrRow{Didáctica de las Ciencias Sociales}{Didáctica General (A); Problemática del Conocimiento Histórico (R); Procesos sociales, políticos, económicos y culturales de la Antigüedad (A); Historia de las Ideas I (A).}
\CorrRow{El mundo y las nuevas territorialidades}{Historia de las Ideas I (R); Procesos sociales, políticos, económicos y culturales de los pueblos originarios de América (R); Procesos sociales, políticos, económicos y culturales de la Antigüedad (R).}
\CorrRow{Procesos sociales, políticos, económicos y culturales del Feudalismo y la Modernidad}{Procesos sociales, políticos, económicos y culturales de la Antigüedad (A); Historia de las Ideas I (A); Procesos sociales, políticos, económicos y culturales de los pueblos originarios de América (R); Problemática del conocimiento histórico (R).}
\CorrTableEnd

\DocQuestion{Tabla 2. Correlatividades de 3\textsuperscript{o} año}
\CorrTableStart
\CorrRow{Hist. de la Educación. Argentina}{Pedagogía (A).}
\CorrRow{Sociología de la Educación}{Pedagogía (R).}
\CorrRow{Práctica Docente III}{Práctica Docente II (A); Psicología Educacional (R); Sujetos de la Educación Secundaria (R); Procesos sociales, políticos, económicos y culturales del Feudalismo y la Modernidad (A); Didáctica de las Ciencias Sociales (R); Procesos sociales, políticos, económicos y culturales Americanos I (R); Historia de las Ideas II (R); El mundo y las nuevas territorialidades (R); Economía Política (R); Todas las UC del Campo específico de 1\textdegree{} año (A).}
\CorrRow{Análisis y Organización de las Instituciones Educativas}{Sujetos de la Educación Secundaria (R); Práctica I (A); Práctica Docente II (R).}
\CorrRow{Procesos sociales, políticos, económicos y culturales Contemporáneos I}{Procesos sociales, políticos, económicos y culturales del Feudalismo y la Modernidad (A); Didáctica de las Ciencias Sociales (A); Procesos sociales, políticos, económicos y culturales Americanos I (R); Historia de las Ideas II (R); Economía Política (R).}
\CorrRow{Procesos sociales, políticos, económicos y culturales Americanos II}{Procesos sociales, políticos, económicos y culturales del Feudalismo y la Modernidad (R); Procesos sociales, políticos, económicos y culturales Americanos I (R).}
\CorrRow{Procesos sociales, políticos, económicos y culturales de Argentina I}{Procesos sociales, políticos, económicos y culturales Americanos I (R); El mundo y las nuevas territorialidades (R).}
\CorrRow{Epistemología de la Historia}{Filosofía (R); Problemática del Conocimiento Histórico (R); Historia de las Ideas II (R).}
\CorrRow{Didáctica de la Historia}{Didáctica de las Ciencias Sociales (A); Sujetos de la educación Secundaria (R).}
\CorrTableEnd

\DocQuestion{Tabla 3. Correlatividades de 4\textsuperscript{o} año}
\CorrTableStart
\CorrRow{Derechos Humanos: Ética y Ciudadanía}{Filosofía (R); Sujetos de la Educación Secundaria (A).}
\CorrRow{Práctica Docente IV}{Todas las Unidades Curriculares de Tercer año. (A).}
\CorrRow{Procesos sociales, políticos, económicos y culturales Contemporáneos II}{Procesos sociales, políticos, económicos y culturales Contemporáneos I (A); Didáctica de la Historia (R); Procesos sociales, políticos, económicos y culturales Americanos II (R); Procesos sociales, políticos, económicos y culturales de Argentina I (R).}
\CorrRow{Procesos sociales, políticos, económicos y culturales Americanos III}{Procesos sociales, políticos, económicos y culturales Contemporáneos I (R); Procesos sociales, políticos, económicos y culturales Americanos II (R); Procesos sociales, políticos, económicos y culturales de Argentina I (R).}
\CorrRow{Procesos sociales, políticos, económicos y culturales de Argentina II}{Procesos sociales, políticos, económicos y culturales Americanos II (R); Procesos sociales, políticos, económicos y culturales de Argentina I (R).}
\CorrRow{Problemáticas Históricas Regionales y Locales}{Procesos sociales, políticos, económicos y culturales Contemporáneos I (R); Didáctica de la Historia (R); Procesos sociales, políticos, económicos y culturales Americanos II (R); Procesos sociales, políticos, económicos y culturales de Argentina I (R); Epistemología de la Historia (R).}
\CorrTableEnd

\DocPara{\textbf{Referencias:} A: aprobada - R: regular.}
\DocPara{\textsuperscript{37} Para cursar una unidad curricular, el estudiante puede tener regularizada/s y/o aprobada/s las UC correlativa/s anterior/es. Para rendir una UC, el estudiante debe tener aprobada/s la/s unidades curriculares correlativas.}
\endgroup
